\subsection{Tangent Planes and Linear Approximations}
\content{

}{% presentation
\vspace{3in}
}{% nopresentation
\vspace{2in}
}{% comments
Derive the equation of the tangent plane. This ends up being 
$$ z-z_0 = f_x(a,b)(x-a)+f_y(a,b)(y-b). $$
}

\content{
\begin{example} Find the tangent plane to the elliptic paraboloid 
$$ z = 2x^2+y^2 $$
at the point $(1,1,3)$.
\end{example} 
}{% presentation
\vspace{3in}
}{% nopresentation
\vspace{2in}
}{% comments
$z-3=4(x-1)+2(y-1)$
}

\content{
The tangent plane can be used to approximate values of a function near a point
whose value is known:
\begin{example} Use the tangent plane to approximate $f(1.1,-0.1)$ where 
$$ f(x,y) = xe^{xy}. $$
\end{example}
}{% presentation
\vspace{3in}
}{% nopresentation
\vspace{2in}
\newpage
}{% comments
approx. $1.1e^{-0.11} \approx 0.98542$.
Write out linearization of $f$ near a point formula!
}

\content{
\begin{example}
Use a tangent plane to approximate the value of $\sqrt{3.05^2 + 3.97^2}$.
\end{example}
}{% presentation
\vspace{3in}
}{% nopresentation
\vspace{2in}
}{% comments
Use function $f(x,y) = \sqrt{x^2+y^2}$.
}

\content{
What happens when the partial derivatives are not continuous at the point
$(a,b)$?
\begin{example}
Let
$$f(x,y) = \begin{cases} \frac{xy}{x^2+y^2} & (x,y) \neq (0,0) \\
0  & (x,y) = 0.
\end{cases} $$ 
Find the first partial derivatives of $f$ at
the point $(0,0)$. What would the linearization of $f$ be at the origin? Is this
a good approximation? 
\end{example}
}{% presentation
\vspace{3in}
}{% nopresentation
\vspace{2in}
}{% comments
}


\content{
Linear appoximations really only make sense for differentiable functions.
\begin{definition}
Let $ \Delta y = f(a+\Delta x) - f(a)$. The function $f(x)$ is said to be differentiable at $a$ if
there exists a function $\epsilon$ such that 
$$ \Delta y = f'(a) \Delta x + \epsilon \Delta x $$
where $\epsilon \rightarrow 0 $ as $\Delta x\rightarrow 0$.
\end{definition}

For functions of two variables, we have 
\begin{definition}
Let $ \Delta z = f(a+\Delta x, b + \Delta y) - f(a,b)$. The function $f(x,y)$ is
said to be differentiable at $(a,b)$ if
there exists a functions $\epsilon_1$ and $\epsilon_2$ such that 
$$ \Delta z = f_x(a,b) \Delta x + f_y(a,b)\Delta y +  \epsilon_1 \Delta x +
\epsilon_2\Delta y $$
where $\epsilon \rightarrow 0 $ as $\Delta x\rightarrow 0$.
\end{definition}
}{% presentation
}{% nopresentation
}{% comments
}

\content{
\begin{theorem}
If $f_x$ and $f_y$ both exist near $(a,b)$ and are continuous at $(a,b)$, then
$f$ is differentiable at the point $(a,b)$.
\end{theorem}
}{% presentation
}{% nopresentation
\newpage
}{% comments
}

